\subsection{Descent: the case of étale spaces}

Let X and Y be topological spaces, and let \(f\colon X\to Y\) be a continuous map. Let T and T' be Y-spaces; endow the X-spaces \(X\times_{Y}T\) and \(X\times_{Y}T'\) with their canonical descent data and denote by \(\mathcal{C}_{f}(X\times_{Y}T;X\times_{Y}T')\) the set of X-morphisms from \(X\times_{Y}T\) into \(X\times_{Y}T'\) compatible with these descent data. For every Y-morphism \(\varphi\colon T\to T'\), the X-morphism \(f^{*}(\varphi)\colon(x,t)\mapsto(x,\varphi(t))\) from \(X\times_{Y}T\) into \(X\times_{Y}T'\) is compatible with the canonical descent data. We will denote by \(f^{*}\colon\mathcal{C}_{Y}(T;T')\to\mathcal{C}_{f}(X\times_{Y}T;X\times_{Y}T')\) the map thus defined.

PROPOSITION 4. --- Suppose that the map f is strict and surjective and that T is an étale Y-space. Then, the map \(f^{*}\) : \(\mathcal{C}_{\mathrm{Y}}(\mathrm{T};\mathrm{T}^{\prime}) \to \mathcal{C}_{f}(\mathrm{X} \times_{\mathrm{Y}} \mathrm{T}; \mathrm{X} \times_{\mathrm{Y}} \mathrm{T}^{\prime})\) is bijective.

Denote by \(\tau\) (resp. \(\tau'\) ) the equivalence relation on \(X\times_{Y}T\) (resp. on \(X\times_{Y}T'\)) which is associated with the canonical descent datum. Since the map f is surjective, the projection \(\mathrm{pr}_2\colon \mathrm{X}\times_{\mathrm{Y}}\mathrm{T} \to \mathrm{T}\) is surjective and the canonical map \((\mathrm{X}\times_{\mathrm{Y}}\mathrm{T})/R_{\tau} \to \mathrm{T}\) is bijective. In particular, the map \(f^*\) is injective. Let us prove that it is surjective. Let \(\varphi\colon \mathrm{X}\times_{\mathrm{Y}}\mathrm{T} \to \mathrm{X}\times_{\mathrm{Y}}\mathrm{T}'\) be an X-morphism compatible with the canonical descent data. For \((x,t) \in \mathrm{X}\times_{\mathrm{Y}}\mathrm{T}\) , we therefore have \(\varphi(x,t) = (x,\overline{\varphi}(t))\) , where \(\overline{\varphi}\) is a map from T to T'.

By definition of \(\overline{\varphi}\) , the map \(\overline{\varphi} \circ \mathrm{pr}_2: \mathrm{X} \times_{\mathrm{Y}} \mathrm{T} \to \mathrm{T}'\) is equal to \(\mathrm{pr}_2 \circ \varphi\) , and is therefore continuous. Since the map f is surjective and strict and T is a Y-étale space, the projection \(\mathrm{pr}_2: \mathrm{X} \times_{\mathrm{Y}} \mathrm{T} \to \mathrm{T}\) is strict (I, p.~32, remark 3). By proposition 9 of I, p.~18, the map \(\overline{\varphi}\) is therefore continuous. It is a Y-morphism such that \(f^*(\overline{\varphi}) = \varphi\) , which proves the proposition.

COROLLARY. --- Let X and Y be topological spaces and let \(f: X \rightarrow Y\) be a strict and surjective map. Let T and T' be Y-étale spaces. If there exists an X-isomorphism \(X \times_{Y} T \rightarrow X \times_{Y} T'\) which is compatible with the canonical descent data, then the Y-spaces T and T' are isomorphic.

Let \(\psi\colon X\times_{Y}T\to X\times_{Y}T'\) be an X-isomorphism of étale spaces. By proposition 4, there exists a unique morphism of Y-spaces \(\varphi\colon T \to T'\) such that \(\psi = f^{*}(\varphi)\) . Since f is surjective, the map \(\varphi\) is bijective. It then follows from cor. 2 of I, p.~30 that \(\varphi\) is an isomorphism, since the Y-spaces T and T' are étale.

PROPOSITION 5. --- Let X and Y be topological spaces and let \(f: X \to Y\) be a continuous map. Suppose that the map f is proper and separated, or open. Every descent datum relative to f on an X-étale space is effective. Moreover, if f is surjective, the quotient space is a Y-étale space.

Let \((Z,p)\) be an X-étale space and let \(\tau\) be a descent datum relative to f on \((Z,p)\) . Denote by \(R_{\tau}\) the equivalence relation in Z associated with \(\tau\) , let \(g: Z \to Z/R_{\tau}\) be the canonical surjection and denote by h the map from Z to \(X \times_{Y} (Z/R_{\tau})\) defined by \(z \mapsto (p(z), g(z))\) . It is continuous and bijective; let us show that it is a homeomorphism.

\begin{enumerate}
\def\labelenumi{\alph{enumi})}
\tightlist
\item
  Suppose first that the map f is open.
\end{enumerate}

Let \(z_0\) be a point of Z; set \(x_0 = p(z_0)\) . Since p is étale, there exists a neighborhood U of \(x_0\) in X and a continuous section \(s: \mathrm{U} \to \mathrm{Z}\) of p above U such that \(s(x_0) = z_0\) . The map \(p \times p: \mathrm{Z} \times_{\mathrm{Y}} \mathrm{Z} \to \mathrm{X} \times_{\mathrm{Y}} \mathrm{X}\) is étale, and the maps from \(\mathrm{U} \times_{\mathrm{Y}} \mathrm{U}\) into \(\mathrm{Z} \times_{\mathrm{Y}} \mathrm{Z}\) defined by \((x, x') \mapsto (s(x), s(x'))\) and \((x, x') \mapsto (s(x), \tau(s(x), x'))\) are continuous sections above \(\mathrm{U} \times_{\mathrm{Y}} \mathrm{U}\) . Since they coincide at every point of \(\mathrm{U} \times_{\mathrm{Y}} \mathrm{U}\) of the form \((x, x)\) , they then coincide on an open neighborhood V of \(\Delta_{\mathrm{U}}\) contained in \(\mathrm{U} \times_{\mathrm{Y}} \mathrm{U}\) . Let \(\mathrm{U}_0\) be a neighborhood of \(x_0\) in X such that \(\mathrm{U}_0 \times_{\mathrm{Y}} \mathrm{U}_0\) is contained in V. Let \((x, u)\) be a point of \(\mathrm{U}_0 \times_{\mathrm{Y}} g(s(\mathrm{U}_0))\) ; let \(x' \in \mathrm{U}_0\) such that \(u = g(s(x'))\) . We therefore have \(s(x) = \tau(s(x'), x)\) , whence \(\mathrm{R}_{\tau}\{s(x), s(x')\}\) and \((x, u) = (x, g(s(x'))) = (x, g(s(x))) = h(s(x))\) .

Since the map g is open (IV, p.~386, prop. 2), the set \(U_{0} \times_{Y} g(s(U_{0}))\) is an open subset of \(X \times_{Y}(Z/R_{\tau})\) on which the map \(h^{-1}\) is equal to the continuous map \(s \circ pr_{1}\) . The map h is therefore a homeomorphism.

\begin{enumerate}
\def\labelenumi{\alph{enumi})}
\setcounter{enumi}{1}
\tightlist
\item
  Now suppose that the map f is proper and separated. Let \(z_{0}\) be a point of Z; set \(x_{0}=p(z_{0})\) and \(y_{0}=f(x_{0})\). The map s given by \(x\mapsto\tau(z_{0},x)\) is a section of p over \(f^{-1}(y_{0})\) . The set \(f^{-1}(y_{0})\) is compact (TG, I, p.~75, th. 1 and I, p.~26, remark 2) and two distinct points of \(f^{-1}(y_{0})\) have disjoint neighborhoods in X, since f is separated (I, p.~25, prop. 1). By Theorem 2 of I, p.~37, there therefore exists a neighborhood \(U_{0}\) of \(f^{-1}(y_{0})\) in X and a section \(s_{0}\) of p over \(U_{0}\) which extends s. The set \(s_{0}(U_{0})\) is an open subset of Z, since p is étale (I, p.~30, cor. 3), and contains the saturation of the set \(\{z_{0}\}\) for the relation \(R_{\tau}\) . The map g is closed (IV, p.~386, prop. 2). There then exists an open set V of \(Z/R_{\tau}\) such that \(W = g^{-1}(V) \cap p^{-1}(U_{0})\) is a neighborhood of \(z_{0}\) contained in \(s_{0}(U_{0})\) (I, p.~75, lemma).
\end{enumerate}

Now let \((x,u)\in\mathrm{U}_{0}\times_{\mathrm{Y}}\mathrm{V}\) and let z be the unique point of Z such that \(h(z)=(x,u)\) ; by definition, one has \(z\in W\) . Since \(\mathrm{W}\subset s_{0}(\mathrm{U}_{0})\) , we have \(z=s_{0}(p(x))\) . This shows that the restriction of \(h^{-1}\) to the open set \(U_{0}\times_{Y}V\) of \(X\times_{Y}(Z/R_{\tau})\) is equal to \(s_{0}\circ p\circ pr_{1}\) . Consequently, h is a homeomorphism.

We have thus shown that the descent datum \(\tau\) is effective. The map f is universally strict (I, p.~20, corollary). Under the hypothesis that f is surjective, it then follows from prop. 8 of I, p.~31 that \(q: Z/R_{\tau} \to Y\) is étale.

